% Generated by report/tools/tables.py from the result files of the repository. Do not edit.
\begin{table}[tbp]
\centering
\footnotesize
\setlength{\tabcolsep}{3.5pt}
\caption[GRPO : vraie récompense contre pile ou face]{Volet GRPO, la comparaison que le protocole demande : gain moyen sous la vraie récompense moins gain moyen sous la récompense tirée à pile ou face, par modèle et par mesure (écart 19 de \file{sim/llm/LLM\_PROTOCOL.md}). Intervalle de Welch à 95~\%, $p$ de la permutation exacte des graines entre les deux récompenses (toutes les partitions), correction de Holm : sur les deux modèles pour pass@1 (contraste principal, $^*$) et pour la justesse des réponses finies, sur les huit contrastes secondaires pour les autres mesures. Établi : intervalle qui exclut zéro et $p$ Holm sous 0,05 ; \textbf{en gras}. Pour Llama, la vraie récompense fait mieux que le hasard parce que la récompense à pile ou face détruit le modèle, pas parce que la vraie récompense le fait progresser (tableau~\ref{tab:grpo-mesures}). Source : \file{sim/llm/results/logs\_stats\_5.json} (\file{sim/llm/logs\_stats.py}).}
\label{tab:grpo-contraste}
\begin{tabular}{@{}lrrrrl@{}}
\toprule
Mesure & \makecell[r]{vraie $-$ hasard\\{}(pts)} & IC de Welch 95~\% & \makecell[r]{$p$ permutation\\{}exacte} & $p$ Holm & Statut \\
\midrule
\multicolumn{6}{@{}l}{\emph{Qwen2.5 0,5B, 5 graines contre 5}} \\
\quad pass@1 (échantillonné)$^*$ & +25 & [−6\,;\,+55] & 0,04 & 0,04 & non établi \\
\quad justesse gloutonne & +24 & [−6\,;\,+54] & 0,05 & 0,12 &  \\
\quad pass@8 & +33 & [−3\,;\,+69] & 0,05 & 0,12 &  \\
\quad part coupée & −64 & [−105\,;\,−23] & 0,02 & 0,10 &  \\
\quad justesse des réponses finies & +27 & [−6\,;\,+59] & 0,04 & 0,04 & non établi \\
\midrule
\multicolumn{6}{@{}l}{\emph{Llama 3.2 1B, 5 graines contre 4}} \\
\textbf{\boldmath \quad pass@1 (échantillonné)$^*$} & \textbf{+30} & \textbf{[+18\,;\,+43]} & \textbf{0,008} & \textbf{0,02} & \textbf{établi} \\
\quad justesse gloutonne & +28 & [+12\,;\,+44] & 0,02 & 0,10 &  \\
\quad pass@8 & +45 & [+32\,;\,+58] & 0,008 & 0,06 &  \\
\quad part coupée & −52 & [−136\,;\,+32] & 0,02 & 0,10 &  \\
\textbf{\quad justesse des réponses finies} & \textbf{+31} & \textbf{[+19\,;\,+43]} & \textbf{0,02} & \textbf{0,04} & \textbf{établi} \\
\bottomrule
\end{tabular}
\end{table}
